\documentclass[9pt, a4paper]{extarticle}
\usepackage[a4paper, top=0in, bottom=0in, left=0.05in, right=0.05in]{geometry} % Minimal margins
\usepackage{amsmath,amssymb,amsthm}
\usepackage{enumitem}
\usepackage{graphicx} % Required for \includegraphics
\usepackage{ragged2e}

\usepackage{fourier}
\usepackage[absolute,overlay]{textpos}
\usepackage{array} % For better column definitions
\usepackage{booktabs} % For improved table aesthetics
\usepackage{subfiles}                % For subfiles compilation

\raggedbottom
\setlist[itemize]{leftmargin=1em}
\newcommand{\dist}[2]{\left\langle #1,\, #2 \right\rangle}
\setlength{\TPHorizModule}{0.1mm} % Set horizontal units to millimeters
\setlength{\TPVertModule}{0.1mm}  % Set vertical units to millimeters
\usepackage{multicol}
\setlength{\parindent}{0pt}
\setlength{\parskip}{5pt plus 1pt minus 1pt}
\usepackage{setspace}
\setstretch{0.98}
\pagestyle{empty}
\setlength{\parskip}{0pt}      % space between paragraphs
\setlength{\parindent}{0pt}    % optional: remove indentation
\setlist{nosep, topsep=0pt, partopsep=0pt, parsep=0pt, itemsep=0pt}
% Redefine the textblock environment
\let\originaltextblock\textblock
\let\endoriginaltextblock\endtextblock
% Adjust vertical centering of the existing \warning symbol
\let\oldwarning\warning
\renewcommand{\warning}{\raisebox{0.3ex}{\oldwarning}}

\makeatletter
\renewcommand{\hline}{\noalign{\vskip 2px}\hrule\@height\arrayrulewidth\@width\linewidth\noalign{\vskip 2px}}
\makeatother

\renewenvironment{textblock}[2][]{%
    \originaltextblock[#1]{#2}%
    \fcolorbox{red}{white}{%
    \begin{minipage}{#2}%
}{%
    \end{minipage}%
    }%
    \endoriginaltextblock
}


\begin{document}
\begin{titlepage}
	\centering
	% Moderately larger size hierarchy
	\makeatletter
	\renewcommand{\Huge}{\@setfontsize\Huge{26pt}{30pt}}
	\renewcommand{\LARGE}{\@setfontsize\LARGE{20pt}{24pt}}
	\renewcommand{\Large}{\@setfontsize\Large{16pt}{20pt}}
	\renewcommand{\large}{\@setfontsize\large{13pt}{16pt}}
	\makeatother

	\vspace*{1cm}
	{\Huge \textbf{Computer Security - CheatSheet}} \\
	\vspace{20pt}
	{\LARGE IN~BA5 - Thomas Bourgeat} \\
	\vspace*{1cm}
	{\Large Notes by Ali EL AZDI} \\
	\vfill

	{\large October 26th, 2025}
	\vspace*{60px}
\end{titlepage}

\small
\noindent
\\
\newpage
\\
\subfile{midterm}\\
\begin{minipage}[t]{0.495\textwidth}
	\textbf{Authentication.} the process of verifying a claimed identity. (Prove who you are.)\\
	\textbf{Passwords.} secret shared between user and system.\\
	\textit{Problems}:
	\begin{itemize}
		\item[-] \textbf{Secure transfer.} password may be eavesdropped when communicated.\\
		      \textit{Encrypting makes password unreadable to eavesdropper, but still allows copy and send.} (\warning \textbf{Replay-Attacks.} sending same request without needing the password)\\
		      \textbf{Challenge-Response} protocol solves this.
		      \begin{enumerate}
			      \item Principal declares they want to login.
			      \item System authenticating principal sends a \textbf{challenge} (a random value).
			      \item Principal encrypts password \textit{together} with the challenge.(ensures password is never encrypted with same value and encryption is fresh)
			      \item After checking password, system authenticating \textit{must} delete the challenge.
		      \end{enumerate}

		\item[-] \textbf{Secure check.} naïve checks may leak information about the password.\\
		      \textbf{Check Options:}
		      \begin{enumerate}
			      \item Letter by letter - Time to check leaks the number of correct checked letters. (Timing channel)
			      \item Check everything - Same time regardless of input (hashing solves this by default).
		      \end{enumerate}

		\item[-] \textbf{Secure storage.} if stolen entire system is compromised.\\
		      Passwords must be stored to be checked against but \textit{not stored in clear} $\rightarrow$ an adversary can read them and impersonate users.\\
		      \textbf{
			      Storage Options :}
		      \begin{enumerate}
			      \item Encrypted Passwords - requires key to be stored in same system as database, not secure.
			      \item Hashed Passwords - Not a solution, hashing does \textbf{not} require a key. \warning passwords are not truly random.$\rightarrow$Database can be downloaded and compared against hash of most common passwords (\textbf{Dictionary Attack)}.
			      \item Hashed + Salted Passwords (Do this !) - \textbf{Salt.} unique random value generated per password. \textbf{We store}: $(\text{username}, H('123456'||\text{salt}_n), \text{salt}_n)$. \\ Dictionary attacks are still possible but same passwords are undetectable.\\
			            \textit{Principal} asks to login$\rightarrow$\textit{System} sends Challenge$\rightarrow$\textit{Principal} sends encrypted password with Challenge $\rightarrow$\textit{System} checks Challenge, decrypts, looks up hash by username, hashes plaintext with salt and checks if same as stored.
		      \end{enumerate}

		\item[-] \textbf{Secure passwords.} easy-to-remember passwords tend to be easy to guess
	\end{itemize}

	\textbf{Biometrics.} measurement and statistical analysis of people's unique physical characteristics(fingerprint, face, \ldots).
	\hline
	\textbf{Biometrics Authentication.}
	\begin{enumerate}
		\item \textbf{Enrollment.} biometric is captured by a sensor and processed to be converted in a stream of bits(Biometric Template). \textbf{Capture} $\rightarrow$ \textbf{Process} $\rightarrow$ \textbf{Store}.
		\item \textbf{Verification.} biometric is captured, processed and compared against the stored one. \textbf{Capture} $\rightarrow$ \textbf{Process} $\rightarrow$ \textbf{Compare}. (threshold for match vs. no match).\\ Balance between False Positives (wrongly authenticated) and False Negatives(wrongly rejected).
	\end{enumerate}
	Storing / Processing can happen locally($\oplus$ privacy) / remotely($\oplus$ security).
	\hline
	\textbf{Authentication Tokens.} by proving ownership of a token. (Smart Card, \ldots)
	\begin{enumerate}
		\item Offline - Device and Server establish a common \textit{seed} (random value) \& synchronize their clocks.
		\item When Device wants identity:
		      \begin{enumerate}
			      \item Server asks to prove it.
			      \item They both compute $n=\frac{{now} - start}{interval}$. \textit{Current interval since start.}
			      \item Device applies a \textbf{keyed} function (pre-shared key) n times on the initial seed.
			      \item Device sends result to Server. (encrypted, eavesdropper can't have the seed).
			      \item Server computes same value and checks against the one received.
		      \end{enumerate}
	\end{enumerate}
	Can't be a hashing function (hash functions aren't keyed!, eavesdropper could compute $h_k(seed)$ for any k after the one observed.).
	\hline

	\begin{minipage}[t]{0.67\textwidth}
		\textbf{Security Engineering process.}
		\begin{enumerate}[leftmargin=10px]
			\item Define a security policy (principals, assets, properties) and a threat model.
			\item Define security mechanisms that support policy of threat model.
			\item Build an implementation supporting the mechanisms
		\end{enumerate}
	\end{minipage}
	\hfill
	\begin{minipage}[t]{0.33\textwidth}
		\textbf{Attack Engineering process.} \\exploit weaknesses introduced during the security
		engineering process due to carelessness, lack of knowledge, or errors.
	\end{minipage}\\
	\textbf{Threat Modelling.} process to identify threats / unprotected resources.
	\begin{itemize}
		\item[-] Attack Trees - represent goal as the root and ways to achieve it as branches.
		\item[-] STRIDE - identify system entities, events, and boundaries of the system. (Spoofing, Tampering, Repudation, Information disclosure, Denial of service, Elevation of privileges).
		\item[-] PASTA - start from buisness goals, processes, and use cases. Find threats in buisness model, assess impact, prioritize on risk.
	\end{itemize}
	\textbf{Methods.}
	\begin{itemize}
		\item[-] \textbf{GET.} requests a resource; parameters are appended to the URL as query key–value pairs after "?", separated by `\&`.
		      \[
			      \underbrace{\texttt{http://}}_{\text{protocol}}
			      \underbrace{\texttt{www.mywebsite.com}}_{\text{hostname / domain}}
			      \underbrace{\texttt{/apparel}}_{\text{directory}}
			      \underbrace{\texttt{/skirt.php}}_{\text{filename}}
			      \underbrace{\texttt{?sku=123\&lang=en\&sect=silk}}_{\text{query parameters}}
		      \]
		\item[-] \textbf{POST.} creates or updates a resource; sends data in the request body (not the URL) and typically modifies server data.\\
		      \texttt{POST /test/demo\_form.php HTTP/1.1}\\
		      \texttt{Host: w3schools.com}\\
		      \texttt{name1=value1\&name2=value2}
	\end{itemize}
	\textbf{Cookies.} small piece of data stored by browser on client (creates HTTP "sessions"). \textit{also used to track users.}. \textbf{Ambient Authority} is used with cookies.
	\textit{logged into bank.com $\rightarrow$ cookies stored for bank.com. \textbf{Any} (!CSRF) new HTTP requests to this website \textbf{will include all cookies} to continue session}.\\
	\textbf{URL.} standardized address that specifies a resource’s access method, host, and path.
	\[
		\underbrace{\texttt{http://}}_{\text{protocol}}\;
		\underbrace{\texttt{www.mywebsite.com}}_{\text{hostname / domain}}\;
		\underbrace{\texttt{/apparel}}_{\text{directory}}\;
		\underbrace{\texttt{/skirt.php}}_{\text{filename}}
	\]
	\textbf{HTML.} a markup language that uses tags to structure content and tell the browser how to render a webpage.\\
\end{minipage}
\begin{minipage}[t]{0.5\textwidth}
	\textbf{PHP.} a server-side scripting language for dynamic web pages; uses variables and request/session data (`\$\_GET`, `\$\_POST`, `\$\_SESSION`) and outputs HTML with `echo`.\\
	Lifetime of a GET/POST request.
	\begin{enumerate}[leftmargin=15px]
		\item Browser sends HTTP GET/POST request.
		\item Web Server receives requests and identifies if it's a static asset(.jpg, etc\dots) or dynamic(.php).
		\item PHP interpreter executes called .php
		\item Sends response to web browser
	\end{enumerate}
	\textbf{Common Weaknesses Enumeration (CWE).}\\ database of software errors leading to vulnerabilities.
	\begin{enumerate}
		\item \textbf{CWE 1 - Insecure Interaction Between Components.} unchecked information sent between different components.
		      \begin{enumerate}
			      \item \textbf{CWE 78 - OS Command Injection.} application embeds unvalidated user input into an OS command, allowing attackers to execute arbitrary commands \textit{ls /home/$user}.

			      \item \textbf{CWE 79 - Cross Site Scripting.} application embeds unvalidated user input into dynamically generated web content, causes browser to execute injected scripts.\\
			            \textbf{Common Forms.}\\
			            \texttt{<img src=x onerror=alert(1)>;}\\
			            \texttt{<svg onload=alert(1)>;}\\
			            \texttt{<body onload=alert(1)>;}\\
			            \texttt{<a href=javascript:alert(1)>click</a>;}
			      \item \textbf{CWE 362 - Cross Site Request Forgery.} malicious website makes a state-changing action using user’s session cookie for another website.\\
			            \fbox{
				            \begin{minipage}[htp]{0.82\textwidth}
					            \textbf{Same Origin Policy (SOP).} web browser security mechanism that restricts scripts from one origin from reading data from another origin. An origin is a (protocol, host, port) \textit{http:\/\/example.com}. \\
					            This doesn't prevent CSRF as it's based on sending a request and that the browser naturally attaches cookies of the destination domain in request (unless special setting).
				            \end{minipage}}\\
			            \textbf{Solutions:}\\
			            For every state-changing endpoint:\\
			            - Check Origin header and only allow exact origin.\\
			            - If Origin missing, require Referer and ensure it starts with expected origin, else reject.\\
			            \textit{Keep GET side-effect free; Require a CSRF token (challenge) on every state-changing request; For high-riks actions require re-auth.}
		      \end{enumerate}
		\item \textbf{CWE 2 - Risky Ressource Management.} not properly managed creation, usage, transfer, or destruction of important system resources

		\item \textbf{CWE 3 - Porous Defenses.} defensive techniques misused, abused, or just ignored
	\end{enumerate}
\end{minipage}
\end{document}

